\section{Conclusion}
\label{sec:conclusion}

% Target: ~half a page. Each TODO below should expand to 2--3 sentences
% that point at concrete numbers from the preceding sections.

\paragraph{Most effective regularization.}
\draftnote{Name the family with the highest test accuracy in
\cref{tab:comparison} and the absolute improvement over baseline. Be
specific about the configuration (e.g. ``RandAugment + Mixup with
$\alpha = 0.2$'').}

\paragraph{Most stable.}
\draftnote{Identify the family whose results varied least across
its hyperparameter sweep (broad plateau in val acc vs hyperparameter)
or, if multi-seed is available, lowest std across seeds. A flat curve
is a feature: it means the practitioner does not need to tune.}

\paragraph{Easiest to implement with high payoff.}
\draftnote{Almost certainly weight decay or label smoothing
(one-line config change). Quantify the payoff per line of code if
possible.}

\paragraph{Which family matters most for ViT.}
\draftnote{Restate the conclusion of \cref{sec:comp:discussion} in
one sentence and connect it to the broader literature: ViT lacks the
inductive biases of CNNs \citep{dosovitskiy2021vit}, so
\emph{data-level} regularization---which effectively injects those
biases through the data distribution---should and does carry the most
weight \citep{touvron2021deit}. \emph{(Verify this claim against your
own results; do not assert it if your results disagree.)}}

\paragraph{Future work.}
\draftnote{1--2 sentences: combining the per-family winners,
extending to longer schedules, or replicating from-scratch training.}
